% =========================================================
% TABLA 2: OBJETIVO ESPECÍFICO 2
% =========================================================
\begin{table}[htbp]
	\centering
	\caption{Cumplimiento del Objetivo Específico 2}
	\label{tab:cumplimiento-objetivo-2}
	\begin{tabular}{|>{\raggedright\arraybackslash}p{0.40\textwidth}|>{\raggedright\arraybackslash}p{0.54\textwidth}|}
		\hline
		\rowcolor[HTML]{1B6B93}
		\multicolumn{2}{|c|}{\color{white}\textbf{\textit{Objetivo específico 2}}}                                                                                                                                                                                                                                                      \\
		\hline
		\multicolumn{2}{|>{\centering\arraybackslash}p{\dimexpr0.94\textwidth+2\tabcolsep+\arrayrulewidth\relax}|}{\textit{Caracterizar alternativas de software libre/código abierto para la implementación de un servicio de directorio en la infraestructura del grupo de investigación GRID}}                                       \\
		\hline
		\textbf{\textit{Detalle de cumplimiento}}                                                                                                                                                                                                                                                     & \textbf{\textit{Observaciones}} \\
		\hline
		Se realizó un Mapeo Sistemático de Literatura (\ac{SMS}) para identificar y caracterizar las alternativas de software libre y código abierto disponibles para la implementación de un servicio de directorio (ver capítulo \textcolor[HTML]{1B6B93}{\textbf{\ref{cap:revision-literatura}}}). &
		\textbf{Paso 1: Revisión sistemática de la literatura.}\par\vspace{0.4em}
		- \textbf{Construcción de la bitácora} y definición de metas, preguntas de investigación y métricas (sección \textcolor[HTML]{1B6B93}{\ref{sec:construccion-bitacora}} y sección \textcolor[HTML]{1B6B93}{\ref{subsec:planeacion}}).\par
		- \textbf{Búsqueda de estudios}: definición de la estrategia, búsqueda en bases de datos y aplicación de criterios de inclusión y exclusión (sección \textcolor[HTML]{1B6B93}{\ref{subsec:estrategia-busqueda}}).\par
		- \textbf{Eliminación de duplicados}, \textbf{priorización de estudios} y \textbf{estrategia de bola de nieve} (sección \textcolor[HTML]{1B6B93}{\ref{sec:eliminacion-duplicados}}, sección \textcolor[HTML]{1B6B93}{\ref{sec:priorizacion-estudios}} y sección \textcolor[HTML]{1B6B93}{\ref{sec:bola-de-nieve}}).\par\vspace{0.6em}
		\textbf{Paso 2: Identificación de tecnologías candidatas.}\par\vspace{0.4em}
		- El \ac{SMS} permitió identificar cuatro alternativas: \textit{OpenLDAP}, \textit{389 Directory Server}, \textit{ApacheDS} y \textit{OpenDJ} (sección \textcolor[HTML]{1B6B93}{\ref{sec:tecnologias}} y sección \textcolor[HTML]{1B6B93}{\ref{subsec:tecnologias-identificadas}}, ver \autoref{tab:tecnologias_candidatas}).\par\vspace{0.6em}
		\textbf{Paso 3: Incorporación de tecnología adicional.}\par\vspace{0.4em}
		- Se incorporó \textit{FreeIPA} como alternativa adicional, al tratarse de una plataforma de código abierto orientada a la gestión centralizada de identidades (sección \textcolor[HTML]{1B6B93}{\ref{sec:incorporacion-tecnologia-adicional}}).\par
		- Se consolidó el conjunto de cinco tecnologías candidatas para su posterior evaluación (sección \textcolor[HTML]{1B6B93}{\ref{sec:conjunto-tecnologias-candidatas}}).                                                                                                                                                          \\
		\hline
	\end{tabular}
\end{table}
